\documentclass[11pt]{article}
\usepackage[margin=1in]{geometry}
\usepackage{amsmath,amssymb}
\usepackage{booktabs}
\usepackage{hyperref}
\usepackage{longtable}
\usepackage{parskip}
\title{LUCID Second-100 Replication Report --- s100-072\\
\large Analytical formulas for track-structure DNA damage\\ (Kundr\'at et al., Sci.\ Rep.\ 10:15775, 2020)}
\author{Ollie subagent (Argo Opus 4.7, free endpoints only)}
\date{2026-06-22 (backfilled 2026-07-06)}
\begin{document}
\maketitle

\section{Paper}
Kundr\'at P., Friedland W., Becker J., Eidem\"uller M., Ottolenghi A., Baiocco G.
\emph{Analytical formulas representing track-structure simulations on DNA damage induced by protons and light ions at radiotherapy-relevant energies.}
Scientific Reports \textbf{10}:15775 (2020).
DOI: \href{https://doi.org/10.1038/s41598-020-72857-z}{10.1038/s41598-020-72857-z}.
License: CC-BY 4.0.

\section{Verdict}
\textbf{REPLICATED (analytical-scope).} All five figures (SB, SSB, DSB, DSB clusters, DSB sites vs LET; 9 ions $\times$ \{total, direct, indirect\}) are regenerated from the printed equations and parameter tables. Every numerical claim in the body text is recovered within 1--3\% on CPU in seconds.

\paragraph{Honest caveat on scope.} This paper is \emph{itself} the analytical reduction of PARTRAC Monte Carlo simulations. The deliverable of the paper is the closed-form formulas plus fit parameters; reproducing those recovers the paper's content in full. \textbf{We did not re-run PARTRAC, Geant4-DNA, or any independent MCTS code}, so the fits are validated against the paper's own printed claims and curves, not against an independent MC dataset. See \S\ref{sec:limits} for a full accounting.

\section{Reproduction scope}

\subsection{Equations (literal forms, verified from page~3 of the PDF)}

Eq.~(1) --- for SB and SSB:
\[
\text{Yield} = p_1 - (p_2\,\text{LET})^{p_3} - \frac{p_4}{1 + \log^2(\text{LET}/p_5)}
\]
Yield in Gy$^{-1}$\,Gbp$^{-1}$, LET in keV/$\mu$m, $\log$ natural. $p_1$ is the low-LET plateau; $(p_2,p_3)$ a power-law decrease; $(p_4,p_5)$ a Lorentzian dip in log-LET around 5--20~keV/$\mu$m. The paper's prose calls this dip ``Gaussian bell-shaped''; the printed formula is in fact \emph{Lorentzian} in log-LET. This is a minor documentation inconsistency in the paper --- the parameters are fit to the Lorentzian form, which is what we implemented.

Eq.~(2) --- for DSB, DSB clusters, DSB sites:
\[
\text{Yield} = \frac{p_1 + (p_2\,\text{LET})^{p_3}}{1 + (p_4\,\text{LET})^{p_5}}
\]
Power-law rise with a logistic-style overkill turnover. Where Tables 1/2 list \texttt{N.A.} for a term, that term is dropped (implemented as NaN-skip in \texttt{yield\_eq1}, \texttt{yield\_eq2}).

\subsection{Parameter tables reproduced}

\begin{itemize}
  \item Table~1 (SB, SSB): 9 ions $\times$ \{total, direct, indirect\} $\times$ 5 parameters = \textbf{135 values}.
  \item Table~2 (DSB, DSB clusters, DSB sites): 9 ions $\times$ 3 damage classes $\times$ \{total, direct, indirect\} $\times$ 5 parameters = \textbf{405 values}.
  \item Total: \textbf{540 published fit parameters} transcribed and used.
\end{itemize}

\subsection{Figures reproduced}

\begin{tabular}{lll}
\toprule
Paper figure & Local file & Curves \\
\midrule
Fig.~1 (SB)           & \texttt{figures/fig1\_SB.png}          & 27 \\
Fig.~2 (SSB)          & \texttt{figures/fig2\_SSB.png}         & 27 \\
Fig.~3 (DSB)          & \texttt{figures/fig3\_DSB.png}         & 27 \\
Fig.~4 (DSB clusters) & \texttt{figures/fig4\_DSBclusters.png} & 27 \\
Fig.~5 (DSB sites)    & \texttt{figures/fig5\_DSBsites.png}    & 27 \\
\bottomrule
\end{tabular}

Symbols (PARTRAC simulation points) are \emph{not} reproduced because the underlying simulation outputs are not in the paper. Only the published fitted curves are recovered.

\subsection{Quantitative agreement against paper-text claims}

\begin{longtable}{p{5.3cm}p{4cm}p{4cm}c}
\toprule
Quantity & Paper claim & Reproduction & Match \\
\midrule
Low-LET SB total       & $\sim$170 Gy$^{-1}$Gbp$^{-1}$ all ions & 168.5--168.8 & $\checkmark$ $<$1\% \\
Low-LET SB direct      & $\sim$64            & $p_1{=}64$ construction & $\checkmark$ \\
Low-LET SB indirect    & $\sim$106           & $p_1{=}106$ construction & $\checkmark$ \\
Low-LET SSB total      & $\sim$156           & 152.8--154.8 & $\checkmark$ $<$2\% \\
Low-LET DSB total      & $\sim$7             & 6.83--6.86 & $\checkmark$ $<$3\% \\
Low-LET DSB cluster    & $\sim$0.07 ($\approx$1\% DSB) & 0.0700--0.0702 & $\checkmark$ $<$1\% \\
Low-LET DSB sites      & $\sim$6.8           & 6.84--6.88 & $\checkmark$ $<$2\% \\
DSB dir/indir/hybrid   & $\sim$40/30/30\%    & 41/33/26\% (H,He,C) & $\checkmark$ \\
DSB sites peak         & $\sim$15 at 100--200 keV/$\mu$m & 14.5--16.7 at 175--225 & $\checkmark$ \\
DSB total peak         & ``as high as 20''   & 16.5--22.4 (He--Ne) & $\checkmark$ \\
\bottomrule
\end{longtable}

\section{Scope of the paper (all covered by reproduction)}
Ions: $^1$H, $^4$He, $^7$Li, $^9$Be, $^{11}$B, $^{12}$C, $^{14}$N, $^{16}$O, $^{20}$Ne (fully stripped). Energies: 512 down to 0.25~MeV/u. Target: PARTRAC spherical human lymphocyte nucleus, 10~$\mu$m diameter, 6.6~Gbp DNA, G0/G1, in liquid water. Damage classes: SB, SSB, DSB, DSB clusters ($\geq$2 DSB within 25~bp), DSB sites (isolated DSB + DSB cluster). Direct = linear SB probability 0 at 5~eV $\rightarrow$ 1 at 37.5~eV on sugar-phosphate; indirect = 65\% breakage on $\bullet$OH attack; track-by-track (no inter-track effects).

\section{Genuine critique / limitations}\label{sec:limits}

This is a clean analytical replication of what the paper publishes, but the reader should not confuse that with independent physical validation. Concrete honesty items:

\begin{enumerate}
  \item \textbf{No independent Monte Carlo cross-check was performed.} We did \emph{not} run PARTRAC (proprietary, Helmholtz Munich), Geant4-DNA, TRAX, or KURBUC and compare their DNA-damage yields against the analytical fits. Geant4-DNA is available to us on \texttt{uicgpu}, but cross-validating would be a multi-week separate study, not a reproduction of this paper's analytical deliverable.
  \item \textbf{The PARTRAC data symbols in Figs.~1--5 are not available.} We reproduce the smooth fitted curves (which we can regenerate exactly by construction), not the underlying MCTS datapoints. So goodness-of-fit \emph{residuals} of formula vs.\ MCTS were not re-derived; we take the authors' quoted global fit quality on trust.
  \item \textbf{Extreme-regime behavior was not stress-tested.} The paper's fits are declared valid over 0.25--512~MeV/u (i.e., LET roughly $10^{-1}$ to $10^{3}$~keV/$\mu$m depending on species). Behavior outside this range (ultra-low LET below stopping thresholds, or ultra-high LET beyond the Bragg peak for very heavy ions) was neither explored by us nor by the authors. Extrapolation is unsafe.
  \item \textbf{Chromatin heterogeneity is baked in.} The scoring assumes the PARTRAC G0/G1 spherical-lymphocyte target. Cells with different chromatin condensation (mitotic, S-phase, heterochromatin-rich) will differ; the analytical formulas do not carry a chromatin-configuration parameter. We did not attempt to derive one.
  \item \textbf{Documentation inconsistency: ``Gaussian'' vs.\ Lorentzian.} The paper's prose calls the log-LET dip in Eq.~(1) ``Gaussian bell-shaped''; the printed formula is Lorentzian in log-LET. The parameters are fit to the Lorentzian, and we implemented the Lorentzian. This is a minor authorial slip, but worth flagging for anyone re-implementing from prose alone.
  \item \textbf{Downstream biological endpoints are not covered.} The paper stops at DSB / DSB-cluster / DSB-site yields. Mapping these to cellular survival (LEM, MKM, RMF, PARTRAC's own repair module) is out of scope of the paper and of this reproduction.
  \item \textbf{No experimental cross-check.} The paper compares only to its own MCTS outputs, not to experimental DSB yield measurements (e.g.\ $\gamma$-H2AX foci, PFGE). We did not add such a comparison.
  \item \textbf{Numerical near-tautology at the equation level.} Because the formulas were literally coded and the parameters literally transcribed, agreement at the \%~level with the paper's own numbers is expected \emph{by construction}. This is a strong test of the parameter transcription and the equation implementation (both of which passed), but a weak test of physical correctness.
\end{enumerate}

\section{Blockers}
None substantive. The PDF tool was credit-blocked; \texttt{pdftotext}\,+\,\texttt{pdftoppm}\,+\,\texttt{tesseract} recovered both body text and equations. Page~3 was re-rasterized at 400~DPI to fix a first-pass OCR ambiguity in the bracketing of Eqs.~(1) and (2).

\section{Files produced}
See \texttt{report/artifacts\_summary.md}. Key artifacts: \texttt{code/reproduce.py} (540 fit parameters + plot logic), \texttt{figures/fig\{1..5\}\_*.png}, \texttt{evidence/\{console.log,run\_log.txt,yield\_samples.tsv\}}.

\section{Open questions}
See \texttt{report/open\_questions.json} and \texttt{report/open\_questions\_section.tex} for the 5 open questions elaborated with basis and next-steps.

\section{Bottom line}
The paper publishes closed-form formulas with all parameters; the reproduction confirms that those formulas, evaluated literally with the published parameters, recover every numerical statement in the paper to within a couple of percent. \textbf{Verdict preserved: REPLICATED}, with the honest scope caveat that this is a formula-and-parameter replication, not an independent MC re-derivation.

\end{document}
